\chapter{Known Results from Literature}\label{ch:literature}

The following simulation edges are known from the literature but not yet
fully formalized in Lean. Edge direction $A \to B$ means ``$A$ simulates $B$''
(equivalently, $B$ embeds into $A$).

\begin{center}
\begin{tabular}{llllll}
\toprule
Edge & $\sovh$ & $\tovh$ & Reference & Status \\
\midrule
TM $\to$ 2-Tag & $\oh{1}$ & bounded & Cocke--Minsky 1964 & hypothesis \\
2-Tag $\to$ CTS & $\times k$ & $2k$ & Cook 2004 & \textbf{proved} \\
CTS $\to$ Rule 110 & $\sim 20$ & $\oh{n}$ & Cook 2004 & unformalized \\
Rule 110 $\to$ CTS & poly & poly & Neary--Woods 2006 & unformalized \\
Rule 110 $\to$ TM(2,3) & $\oh{w}$ & $\oh{w}$ & Smith 2007 & hypothesis \\
TM $\leftrightarrow$ GS & $1$ & $1$ / $\oh{w}$ & Moore 1991 & \textbf{proved} / 1 sorry \\
ECA 110 $\leftrightarrow$ ECA 124 & $1$ & $1$ & mirror & \textbf{proved} \\
Game of Life $\to$ TM & $\sim 10^3$ & $\oh{1}$ & Rendell 2002 & unformalized \\
Billiard $\to$ TM & $\oh{1}$ & $\oh{|Q|}$ & Miranda--Ramos 2025 & unformalized \\
\bottomrule
\end{tabular}
\end{center}

The Neary--Woods polynomial chain:
$\mathrm{TM} \xrightarrow{\oh{1}} \mathrm{ClockwiseTM} \xrightarrow{\oh{n}}
\mathrm{CyclicTag} \xrightarrow{\oh{n}} \mathrm{Rule\;110}$,
achieving total polynomial overhead $\oh{t^2}$, replacing the exponential
Cocke--Minsky path. This demonstrates \emph{edge discovery} as a central
graph operation: finding a better path changes the overhead landscape.
